\begin{abstract}
Deep learning models for visual recognition often exhibit systematic errors due to underrepresented semantic subpopulations. While existing debugging frameworks can identify these failure slices, effectively repairing them remains difficult. Current solutions often rely on manually designed prompts to generate synthetic images—an approach that introduces distribution shift and semantic errors, often creating new bugs. To address these issues, we introduce \model, a framework for distribution-consistent model repair via controlled  generation. Our approach employs a structural-conditioned diffusion model to generate semantically faithful images that modify only specific failure attributes while preserving the underlying data distribution. To ensure the reliability of the repair data, we implement a  verification mechanism using a large vision--language model (LVLM) to enforce semantic consistency and label preservation. By retraining models on the synthetic data, we significantly reduce errors in rare cases and improve overall performance. Our experiments show that \model achieves superior robustness by maintaining high precision in attribute editing without introducing additional bugs.
\end{abstract}

\section{Introduction}
Despite achieving high accuracy on standard benchmarks, computer vision models often fail on rare semantic subpopulations~\cite{barbu2019objectnet,gao2023adaptive,leclerc20223db}. These failures usually stem from dataset bias, where specific attribute combinations (e.g., darker-skinned individuals with specific hair textures) are underrepresented in the training set~\cite{buolamwini2018gender}. Identifying and repairing these "bug slices" is essential for the reliable deployment of AI systems in safety-critical applications.

Recent research has focused on identifying these failures using interpretable debugging pipelines~\cite{gao2023adaptive,chen2023hibug,singla2024data}. For instance, \hibug~\cite{chen2023hibug} uses vision--language models to discover failure slices based on shared visual attributes. However, the subsequent repair process remains a major bottleneck. Existing repair strategies generally fall into two categories: retrieval-based and generation-based. Retrieval-based methods~\cite{gao2023adaptive,singla2024data} collect samples from external datasets, which often introduces a domain gap. Generation-based methods, such as those suggested by \hibug, use text-to-image models to synthesize training data. However, standard generative augmentation suffers from two critical flaws: (1) \textbf{Semantic Drift}: Generative models often fail to preserve attributes unrelated to the intended edit, and (2) \textbf{Distribution Mismatch}: Synthetic images may not align with the original data distribution, leading the model to learn "generative artifacts" rather than the intended semantic concepts.

To overcome these challenges, we propose \model, we propose \model, a targeted model repair method designed to correct failures caused by underrepresented semantic subpopulations through controlled image generation—without degrading overall performance or introducing new errors. Compared to AdaVision~\cite{gao2023adaptive} and DCD~\cite{singla2024data}, which retrieve samples from another dataset, our method generates new examples grounded in the training set to ensure distributional consistency. Unlike \hibug~\cite{chen2023hibug}, which simply composes failure attributes into language prompts, \model anchors generation on real instances in the training set to preserve unrelated attributes and avoid unintended changes. While the latent-space-based filtering method~\cite{jain2023latentfailures} explores repair via diffusion models, their approach relies solely on caption-based prompting and does not guarantee the semantic accuracy of critical attributes. In contrast, \model mitigates the unreliability of generative models in rendering sensitive attributes—such as skin tone or hair color—by further filtering out semantically inconsistent images using a large vision--language model (LVLM)~\cite{bai2025qwen25vl, liu2023improvedllava}, and the LVLM is verified by human audit. This  system ensures that the final outputs faithfully reflect the intended attribute change (e.g., generating a sad expression on a darker-skinned woman with red hair, whereas standard diffusion models often mistakenly produce a darker-skinned man with black hair instead), enabling precise and robust repair of underrepresented semantic subpopulation failures.

\model directly addresses two major challenges in existing model repair strategies: {\bf 1)} ensuring the generated images reflect the true failure attributes, and {\bf 2)} aligning them with the original data distribution. By retraining models on this augmentation set, we reduce errors associated with underrepresented semantic subpopulations, where model failures stem from insufficient training coverage of specific attributes. The contributions of this work are as follows:
\begin{itemize}
\item We formulate a targeted model repair pipeline, \model, which leverages conditional text-to-image generation and LVLM-based filtering to synthesize high-quality, attribute-faithful data for correcting model failures arising from underrepresented subpopulations.

\item We introduce a  verification mechanism that leverages LVLMs to mitigate the inherent unreliability of generative models, ensuring that the synthesized repair data is both semantically accurate and remains aligned with the original dataset distribution.

\item We demonstrate that our approach significantly enhances model performance on rare-case failure slices across multiple architectures and datasets, outperforming current baselines in both repair effectiveness and the prevention of new bugs.
\end{itemize}